\documentclass[11pt,a4paper]{article}

\usepackage{kotex}
\usepackage{fontspec}
\setmainfont{Times New Roman}
\setmainhangulfont{AppleMyungjo}
\setsanshangulfont{Apple SD Gothic Neo}

\usepackage[a4paper,top=18mm,bottom=18mm,left=20mm,right=20mm]{geometry}
\usepackage{fancyhdr}
\setlength{\headheight}{24pt}
\usepackage{enumitem}
\usepackage{mdframed}
\usepackage{array}
\usepackage{booktabs}

\setlength{\parindent}{1em}
\linespread{1.18}

\pagestyle{fancy}
\fancyhf{}
\renewcommand{\headrulewidth}{0.6pt}
\fancyhead[L]{\sffamily\bfseries 2026 고1 영어 변형 — 정답 및 해설 (지문 32)}
\fancyhead[R]{\sffamily [다양한 장소에서 공부하기]}
\renewcommand{\footrulewidth}{0pt}
\fancyfoot[C]{\framebox[16mm][c]{\thepage}}

\newmdenv[linewidth=0.6pt,roundcorner=3pt,innertopmargin=7pt,innerbottommargin=7pt,
          innerleftmargin=9pt,innerrightmargin=9pt,skipabove=5pt,skipbelow=5pt]{pbox}

\setlength{\parindent}{0pt}

\begin{document}

\begin{center}
{\fontsize{15}{17}\selectfont\sffamily\bfseries [지문 32] 정답 및 해설}
\end{center}
\vspace{6pt}

%% ── 정답 표 ──────────────────────────────────────────
\begin{center}
\begin{tabular}{>{\centering\arraybackslash}p{18mm}
                >{\centering\arraybackslash}p{28mm}
                >{\centering\arraybackslash}p{18mm}
                >{\centering\arraybackslash}p{18mm}}
\toprule
문항 & 유형 & 정답 & 배점 \\
\midrule
1 & 빈칸추론 & ④ & 2점 \\
2 & 영어 제목 & ② & 2점 \\
3 & 서술형 — 해석 & (아래 참조) & 2점 \\
4 & 내용 불일치 & ④ & 2점 \\
\bottomrule
\end{tabular}
\end{center}

\bigskip

%% ── 문항별 해설 ──────────────────────────────────────
{\sffamily\bfseries 1번 해설 (빈칸추론, 정답 ④ consolidate)}

\medskip
\begin{pbox}
\textbf{정답 근거}: ``that rich experience can \underline{\hspace{2cm}} what you are learning''에서
빈칸은 `다양한 장소에서 공부하는 풍부한 경험'이 학습에 미치는 효과를 나타낸다.
지문은 실험 결과를 통해 장소를 바꾸어 공부한 학생들이 단어를 훨씬 더 잘 기억했다고 서술하고
있으므로, 빈칸에는 ``학습한 내용을 강화·공고히 하다''는 의미의 \textbf{④ consolidate}가 적절하다.
`reinforce'와 동의어로, 기억 정착(memory retention)을 강화하는 학문적 표현이다.
\end{pbox}

\medskip
\textbf{오답 소거}
\begin{itemize}[leftmargin=2em,itemsep=1pt]
  \item ① replace — `대체하다': 경험이 학습 내용을 대체한다는 것은 문맥에 어긋난다.
  \item ② complicate — `복잡하게 하다': 지문의 긍정적 효과와 반대 방향이다.
  \item ③ contradict — `모순되다': 경험이 학습을 부정한다는 의미로 전혀 맞지 않는다.
  \item ⑤ accelerate — `가속화하다': 속도보다 깊이·보유(retention)가 핵심이므로 부정확하다.
\end{itemize}

\bigskip

{\sffamily\bfseries 2번 해설 (영어 제목, 정답 ②)}

\medskip
\begin{pbox}
\textbf{정답 근거}: 지문의 핵심은 ``여러 장소에서 공부하면 더 잘 기억한다''이다.
\textbf{② Study in Multiple Places to Remember More}는 이 주제를 가장 직접적으로 표현한다.
\end{pbox}

\medskip
\textbf{오답 소거}
\begin{itemize}[leftmargin=2em,itemsep=1pt]
  \item ① `편안한 환경이 집중력을 높인다' — 지문은 오히려 같은 장소를 고집하지 말라고 한다.
  \item ③ `왜 학생들이 잊어버리는가' — 망각의 원인이 아니라 기억력 향상법이 주제이다.
  \item ④ `그룹 학습이 혼자 학습보다 낫다' — 집단 학습은 지문에서 전혀 언급되지 않는다.
  \item ⑤ `전용 공부 공간 만들기' — 지문의 메시지(장소를 바꿔라)와 정반대이다.
\end{itemize}

\bigskip

{\sffamily\bfseries 3번 해설 (서술형 — 해석)}

\medskip
\textbf{대상 문장}

\begin{pbox}
\itshape
Variety creates rich association, even when those connections form in the
background, totally outside of what we are consciously thinking.
\end{pbox}

\medskip
\textbf{모범 답안}

\begin{pbox}
다양성은 풍부한 연상을 만들어 내는데, 심지어 그 연결들이 우리가 의식적으로 생각하는 것의
완전히 바깥인 배경에서 형성될 때에도 그러하다.
\end{pbox}

\medskip
\textbf{채점 포인트} (각 1점씩, 총 2점)
\begin{itemize}[leftmargin=2em,itemsep=1pt]
  \item \textbf{[1점]} ``creates rich association'' → `풍부한 연상을 만들어 낸다(생성한다)' 정확히 표현
  \item \textbf{[1점]} ``totally outside of what we are consciously thinking'' →
        `우리가 의식적으로 생각하는 것의 완전히 바깥에서' 정확히 표현
        (\,`배경에서'는 bonus; `consciously' 누락 시 감점)
\end{itemize}

\vspace{8pt}
\noindent{\sffamily\bfseries 4. 내용 불일치 - 정답: ④}
\par 본문은 한 장소에 계속 머무르기보다 여러 장소에서 공부하는 것이 기억에 도움이 된다고 말한다. 따라서 항상 좋아하는 한 구석으로 돌아가라는 ④는 본문과 반대이다.

\end{document}
